\documentclass[11pt,letterpaper]{article}
\usepackage[T1]{fontenc}
\usepackage[utf8]{inputenc}
\usepackage{lmodern}
\usepackage[margin=1in]{geometry}
\usepackage{microtype}
\usepackage{xcolor}
\usepackage{hyperref}
\hypersetup{
  colorlinks=true,
  linkcolor=blue,
  urlcolor=blue,
  citecolor=blue,
  pdftitle={Plale Lab Shared Glossary and Operational Definitions},
  pdfauthor={},
  pdfsubject={Working proposal for team review; not lab-approved}
}
\setlength{\parindent}{0pt}
\setlength{\parskip}{0.5em}
\newcommand{\entryfield}[2]{\par\textbf{#1.} #2\par}
\title{Plale Lab Shared Glossary\\and Operational Definitions}
\author{Working proposal for team review}
\date{Starter version 0.1 --- 10 October 2026}

\begin{document}
\maketitle
\begin{center}
\fbox{\parbox{0.9\linewidth}{\textbf{Status: proposed; not lab-approved.}
This starter offers a format and illustrative definitions for discussion.
It does not claim team consensus or establish lab policy.}}
\end{center}

\section{Purpose and use}
A shared glossary can make the intended meaning and boundaries of a term
visible. An operational definition adds the observations or decision rules
needed to apply that meaning in a specified setting. Neither a concise
definition nor a successful document build establishes scientific validity.

The entries below are starter proposals. Project-specific definitions,
thresholds, and examples should be added through review, with public sources
and a stated scope. Where evidence or a decision is missing, record that gap
explicitly rather than treating it as settled.

\section{Proposed entry and review convention}
Each entry records the term, working definition, scope, inclusions,
exclusions, operational criteria where applicable, source and source version,
entry version, and status. An entry may also record a synthetic example,
counterexample, related terms, and open questions.

New entries start as \textbf{proposed}. An \textbf{adopted} entry requires an
explicit team decision with a date and public reference. Adoption is limited
to the entry's stated scope. A merged pull request or a passing PDF build is
not, on its own, a decision to adopt a definition. This review convention is
itself a proposal for team consideration.

Copy \texttt{templates/definition-entry.tex} into the appropriate part of
this document, replace the placeholders, and cite public sources through
\texttt{references.bib}. Keep private material outside this public repository.

\newpage
\section{Illustrative proposed entries}
\subsection{Working definition}
\label{term:working-definition}
\entryfield{Working definition}{An explicit, revisable statement of the
meaning intended for a term within a stated context.}
\entryfield{Scope}{Terminology entries in this glossary.}
\entryfield{Includes}{The intended referent, relevant context, and boundaries
needed to interpret use of the term.}
\entryfield{Excludes}{Claims that the meaning is universal, settled across
disciplines, or already adopted by the lab.}
\entryfield{Operational criteria}{For this glossary, a candidate entry states
its meaning and scope and identifies inclusions and exclusions. If a term
will be used to measure or classify observations, the entry also supplies
the operational criteria or explicitly marks them as pending. This is a
proposed editorial check, not a scientific measurement rule.}
\entryfield{Source and version}{Original starter proposal in this repository,
version 0.1, 10 October 2026; no external authority is claimed.}
\entryfield{Entry version}{0.1.}
\entryfield{Status}{Proposed. Adoption decision: none recorded.}

\subsection{Operational definition}
\label{term:operational-definition}
\entryfield{Working definition}{A specification of how a concept will be
observed, measured, or classified in a stated setting, including the rule
used to interpret the observations.}
\entryfield{Scope}{Definitions used to support a repeatable measurement or
classification procedure.}
\entryfield{Includes}{The observation unit, inputs, procedure, decision rule,
and treatment of ambiguous or missing information, where relevant.}
\entryfield{Excludes}{A broad description of a concept without a procedure;
a claim that a procedure alone establishes the concept's validity.}
\entryfield{Operational criteria}{A proposed procedure identifies the unit
and inputs, states executable steps, specifies thresholds and units when
used, and explains ambiguous and missing cases. It includes a synthetic
example and counterexample when useful. Any unsettled rule is marked pending.}
\entryfield{Source and version}{Original starter proposal in this repository,
version 0.1, 10 October 2026; no external authority is claimed.}
\entryfield{Entry version}{0.1.}
\entryfield{Status}{Proposed. Adoption decision: none recorded.}

\newpage
\subsection{Provenance}
\label{term:provenance}
\entryfield{Working definition}{Information about the origin and production
history of an entity, including relevant activities and responsible agents.}
\entryfield{Scope}{Describing the derivation of a public document, definition,
or other research object. This entry is a proposed local interpretation
informed by the W3C PROV-O Recommendation \cite{w3c-provo}; it is not a
claim of conformance to the full standard.}
\entryfield{Includes}{Source identifiers and versions, derivation steps, and
attribution where publicly shareable.}
\entryfield{Excludes}{An unsupported origin claim; a guarantee of accuracy
or trustworthiness; disclosure of restricted source material.}
\entryfield{Operational criteria}{For a glossary entry, record a public source
or identify an original proposal; record the source version or date; and
describe how the proposed definition was adapted. Use a source locator
such as a section, DOI, or stable URL where available. Mark gaps as unknown.}
\entryfield{Source and version}{W3C, \emph{PROV-O: The PROV Ontology},
Recommendation of 30 April 2013, Sections 1--3 \cite{w3c-provo};
entry adaptation version 0.1, 10 October 2026.}
\entryfield{Entry version}{0.1.}
\entryfield{Status}{Proposed. Adoption decision: none recorded.}

\section{Questions for team review}
\begin{itemize}
\item Which terms are most useful to define first, and for which projects?
\item What decision process should establish adoption and later revision?
\item Which entries need measurable operational criteria or boundary examples?
\end{itemize}

\bibliographystyle{plain}
\bibliography{references}
\end{document}
